% Pursuing Stacks (cotes 134-2 to 134-8): archive cote 134-2, site page 63 (PDF page 64).
% appendix (letters to Breen). Draft compared: ps1.tex:1468-1474.
% AI-transcribed by Claude Opus 5.5 (claude-opus-5-5, adaptive thinking, effort high), Message Batches API, batch ps1 (msgbatch_01Rv5hCaNTtUt4JHx1yqoHPL), 2026-11-12; not yet checked by eye.
\scanpage{134-2/63}
\mspage{-18 -}

\margin{G'} such a group, then $G^{\over{\unsure{/}}{\unsure{\prime}}} = \mathrm{Ext}^1(G, \mathbb{Q}/\mathbb{Z})$ can canonically be given a structure
of quasi-algebraic group (i.e. defined mod radical isogeny), doubtless in
a unique manner provided it verifies some functorial properties, and on
requiring that for $G = \mathbb{G}_a$ one finds that $\mathrm{Ext}^1(\mathbb{G}_a, \mathbb{Q}/\mathbb{Z}) \simeq \mathbb{G}_a$ with the
\margin{G'} usual structure. Let $\Delta G = G^{\over{\unsure{/}}{\unsure{\prime}}} = \mathcal{E}xt^1(G, \mathbb{Q}/\mathbb{Z})$. One finds $G \simeq \Delta\Delta G$ i.e. $\Delta$
is an authentic autoduality! I call $\Delta$ \textit{Serre duality}. It surely goes over
to ind-progroups on an arbitrary base field (not necessarily algebraically closed)
in the case $p > 0$ . Moreover, for finite étale groups, it is $\mathcal{E}xt^0(G, \mathbb{Q}/\mathbb{Z})$
(Pontrjagin duality) which gives a perfect duality. One could \deleted{assemble}\inserted{screw together}, in
an appropriate derived category, Serre duality and Pontrjagin duality, by
taking $G \longmapsto \Delta G = \mathbb{R}\mathcal{H}om\,(G, Q/\mathbb{Z})$: one calls this (``cohomological'') Serre
duality. This will be a magnificent autoduality, if one puts oneself in a
derived category where the $\underline{H}^i$ of the envisaged complexes are (up to
passing to the limit) extensions of étale groups by connected unipotent
groups. Now one \deleted{\unsure{only meets}} \inserted{gets only} such complexes, \over{\unsure{on}}{by} ``integrating'' finite coefficients
$F$ on $X$ by $Rf_!$ or $Rf_*$ . This being said, by \deleted{taking \unsure{i}} \inserted{passing} to the limit in
the initial formulation (or equivalently by replacing the $(\mathbb{Z}/n\mathbb{Z})_k$ ,
previously considered, by $(\mathbb{Q}/\mathbb{Z})_k$ on $k$ , and forming $f^!(\mathbb{Q}/\mathbb{Z})_k = (\mu_{\hand{\unsure{\infty}}})_X$)
the duality formula takes the form
\[ \Delta(Rf_!(F)) \simeq Rf_*(DF[2]) \qquad \hand{$\curvearrowleft$}\text{``shift'' of dimension} \]
where $D$ is the ``Cartier duality'' $\uline{R\mathcal{H}om}(F, \mu_{\hand{\unsure{\infty}}})$ (or $R\mathcal{H}om(F, \mathbb{G}_m)$ if one
prefers?), and $\Delta$ is the Serre duality: cohomology with proper supports and \hand{with}
arbitrary supports are exchanged by duality, when one takes upstairs Cartier
duality, and downstairs Serre duality.

The validity of the duality formula is not open to doubt - the principal
work for establishing it consists certainly in a careful description of \hand{the}
\deleted{coefficient} categor\over{y}{ies} \inserted{of coefficients} with which one is working, as well on $X$ as on $k$ ,
and of the functors $D$ and $\Delta$ . As the definition of an arrow is immediate,
once the building of the machine has been accomplished, the validity of the

\margin{63}
